\section{Mechanism Classes and the Screening Protocol}\label{sec:signatures}
% Table 2: mechanism classes and their physical signatures (Plan §9; plan/CHANGELOG.md C6).
% Cell sources: paper/tables/table2_NOTES.md. Needs \usepackage{booktabs,tabularx} in main.tex.
\begin{table}[t]
\caption{Predicted response of each mechanism class to each probe of a frozen deadlock. ``---'': the probe does not test that class.}
\label{tab:signatures}
\centering
\scriptsize
\setlength{\tabcolsep}{2pt}
\renewcommand{\arraystretch}{1.15}
\begin{tabularx}{\textwidth}{@{}l*{7}{>{\raggedright\arraybackslash}X}@{}}
\toprule
Class
 & Quench
 & Replay, $10^{-9}$\,m
 & $\Delta t$, per-step noise
 & $\Delta t$, $\sqrt{\Delta t}$ noise
 & Rest-drive push, $+$ vs $-$
 & Tilt $b\to-b$
 & Corrected numerics \\
\midrule
P
 & held
 & held
 & held
 & held
 & unchanged
 & unchanged
 & persists \\
L
 & frees, deterministic
 & identical replays
 & no $\Delta t$ dependence
 & no $\Delta t$ dependence
 & $+$ frees, $-$ does not
 & ---
 & persists \\
T
 & frees at random, exponential times
 & replays disagree
 & rate falls as $\Delta t$ grows
 & no $\Delta t$ dependence
 & ---
 & ---
 & persists \\
D
 & frees only beyond $b_f$
 & identical replays
 & no $\Delta t$ dependence
 & no $\Delta t$ dependence
 & ---
 & first out switches side
 & persists \\
N
 & frees or holds, depending on solver details
 & outcome changes under $10^{-9}$\,m
 & ---
 & ---
 & ---
 & ---
 & effect disappears \\
\bottomrule
\end{tabularx}
\end{table}

%
By which mechanism does a perturbation free a frozen deadlock, and how are two perturbations with the same rate told apart?

\subsection{Mechanism Classes}
A deadlock snapshot is a rest state of Eq.~\eqref{eq:force}: every agent is still and the forces on it balance. At a lone keystone, removing the blocker frees the others; at a second trap another agent takes its place and is held; in a group standoff several agents stay held. A perturbation acts on a deadlock in one of five ways (responses to each probe in Table~\ref{tab:signatures}). \textbf{P} (prevent): it acts only while the crowd moves and so changes only how often deadlocks form. \textbf{L} (reshape): a deterministic force that persists at rest moves or removes the fixed point. \textbf{T} (heat): added fluctuations leave the fixed point in place and carry the state over a barrier $\Delta U$ at a rate set by $\Delta U/k_BT_{\rm eff}$. \textbf{D} (steer): a finite directional bias breaks a symmetry and removes a strict minimum at a finite threshold $b_f$. \textbf{N} (glitch): the outcome changes under a perturbation far below the model's physical scales, or the effect vanishes in the corrected implementation.
The Jaccard index $J$ of the outcome sets (Sect.~\ref{sec:stats}), the time to free, the final-state class and stall survival are descriptive signatures beyond Table~\ref{tab:signatures}; $J$ separates two perturbations only at the pair standoff, where both are deterministic.

\subsection{The Screening Protocol}\label{sec:protocol}
The protocol screens a resolver, the steering mechanism credited with resolving deadlocks, in five steps (Fig.~\ref{fig:protocol}).
\begin{enumerate}
\item Build a snapshot library of frozen deadlocks (Sect.~\ref{sec:model}).
\item Apply the probes to each perturbation on the same snapshots: the quench, the $10^{-9}$\,m replay, a change of $\Delta t$ with per-step and with $\sqrt{\Delta t}$-scaled noise, the rest-drive push $\pm$, the lateral tilt $\pm b$ and the corrected implementation (here the RVO2 solver). The rest-drive push is a constant force equal to ORCA's rest drive, towards the door ($+$) or away from it ($-$).
\item Read each perturbation's signatures, its responses to the probes.
\item Assign each perturbation the class whose signatures it shows (Table~\ref{tab:signatures}).
\item Compare the resolver with its rate-matched surrogate, a perturbation of another class whose single parameter is tuned to the resolver's rate on the same deadlocks (freed fraction for snapshots, completion fraction for whole runs). The two share a mechanism only if no probe separates them. A probe separates two perturbations when their 95\% intervals for that signature exclude each other, or when one shows the predicted effect and the other its absence; replay and handedness are exact counts.
\end{enumerate}
The rate match only calibrates the surrogate. The rest-drive push is the reference perturbation of known class L (Sect.~\ref{sec:att-lone}).

\begin{figure*}[t]\centering% Fig. (protocol): the five steps of the screening protocol as a greyscale vector flowchart, full \textwidth.
% Input from paper/sections/04_protocol.tex inside a figure* environment; needs tikz with the libraries
% arrows.meta and calc (both loaded in paper/main.tex). No fills: black text and arrows, grey frames.
% One row per step (step name left, detail right); the rows carry the detail of the protocol so that the figure
% stands alone without a numbered list in the text.
\begin{tikzpicture}[
    font=\scriptsize,
    every node/.style={outer sep=0pt},
    frame/.style={draw=black!60, line width=0.5pt},
    rule/.style={draw=black!60, line width=0.4pt},
    flow/.style={-{Stealth[length=1.5mm,width=1.2mm]}, line width=0.6pt, color=black},
  ]
  \pgfmathsetlengthmacro{\hw}{0.17\textwidth}%       heading cell
  \pgfmathsetlengthmacro{\pad}{3pt}%
  \pgfmathsetlengthmacro{\bw}{\textwidth-\hw-2*\pad-0.5pt}% body text width
  \pgfmathsetlengthmacro{\vgap}{0.30cm}%
  \pgfmathsetlengthmacro{\hhw}{0.5*\hw}%
  % \protostep{name}{north-west corner of the body cell}{heading}{body}
  \long\def\protostep#1#2#3#4{%
    \node[anchor=north west, text width=\bw, inner xsep=\pad, inner ysep=2.6pt, align=left, execute at begin node={\hyphenpenalty=10000 \exhyphenpenalty=10000 \relax}] (#1b) at #2 {#4};%
    \draw[frame] ($(#1b.north west)-(\hw,0)$) rectangle (#1b.south east);%
    \draw[rule] (#1b.north west) -- (#1b.south west);%
    \node[anchor=west, inner sep=0pt, font=\scriptsize\bfseries] at ($(#1b.west)-(\hw,0)+(\pad,0)$) {#3};%
    \coordinate (#1t) at ($(#1b.north west)-(\hhw,0)$);%
    \coordinate (#1s) at ($(#1b.south west)-(\hhw,0)$);%
  }%
  \protostep{s1}{(\hw,0)}{1\hspace{0.5em}Snapshots}%
    {A library of frozen deadlocks, each a rest state saved from a whole run: lone keystones, second traps and pair standoffs.}
  \protostep{s2}{($(s1b.south west)-(0,\vgap)$)}{2\hspace{0.5em}Probes}%
    {Each perturbation is applied to the same snapshots under each probe: quench (restart under the perturbation); replay (a copy displaced by $10^{-9}$\,m); time step $\Delta t$ changed, with per-step and with $\sqrt{\Delta t}$-scaled noise; rest-drive push $\pm$ (a constant force equal to the rest drive of force-coupled ORCA, towards ($+$) or away from ($-$) the door); lateral tilt $\pm b$; corrected implementation (the RVO2 solver).}
  \protostep{s3}{($(s2b.south west)-(0,\vgap)$)}{3\hspace{0.5em}Signatures}%
    {The responses of each perturbation to the probes, paired by snapshot: freed fraction at horizon $T$, times to free, outcome changes under replay and handedness under tilt.}
  \protostep{s4}{($(s3b.south west)-(0,\vgap)$)}{4\hspace{0.5em}Class}%
    {The class whose signatures the perturbation shows (Table~\ref{tab:signatures}): Prevent; Reshape, axial or lateral; Heat; Glitch.}
  \protostep{s5}{($(s4b.south west)-(0,\vgap)$)}{5\hspace{0.5em}Compare}%
    {The resolver (the mechanism under test) against its surrogate, a perturbation of another class whose one parameter is set so that it is fraction-matched at $T$ on the same deadlocks. The two share a mechanism only if no probe separates them (paired tests, Holm-corrected across probes).}
  \foreach \a/\b in {s1/s2, s2/s3, s3/s4, s4/s5}{\draw[flow] (\a s) -- (\b t);}%
\end{tikzpicture}%
\caption{Screening protocol: the mechanism under test and its rate-matched surrogate are probed on the same frozen deadlocks and each is assigned a class.}\label{fig:protocol}\end{figure*}

\subsection{Terms That Vanish at Rest}
Stack the $n$ residual agents into $\mathbf x,\mathbf v\in\mathbb R^{2n}$, freeze the densities at their snapshot values and split the force as
\begin{equation}\label{eq:split}
\mathbf F(\mathbf x,\mathbf v)=\mathbf F_0(\mathbf x)-\mathsf D\,\mathbf v+\mathbf G(\mathbf x,\mathbf v),
\end{equation}
with the rest field $\mathbf F_0(\mathbf x)=\mathbf F(\mathbf x,\mathbf 0)$, the damping $\mathsf D={\rm diag}(m_i/\tau_i)$ of Eq.~\eqref{eq:des} and $\mathbf G(\mathbf x,\mathbf 0)=\mathbf 0$.

\begin{proposition}[rest-state lemma]\label{prop:rest}
Suppose $\mathbf v\cdot\mathbf G(\mathbf x,\mathbf v)\le0$ for all $|\mathbf v|<\epsilon$. Then (i) the fixed points of $\mathsf M\dot{\mathbf v}=\mathbf F$, $\dot{\mathbf x}=\mathbf v$, with $\mathsf M$ the mass matrix, are exactly the states $(\mathbf x^\ast,\mathbf 0)$ with $\mathbf F_0(\mathbf x^\ast)=\mathbf 0$; and (ii) if $\mathbf F_0=-\nabla U$, every strict local minimum of $U$ is a Lyapunov-stable fixed point.
\end{proposition}
\begin{proof}
(i) A fixed point has $\mathbf v=\mathbf 0$ and $\mathbf F(\mathbf x^\ast,\mathbf 0)=\mathbf F_0(\mathbf x^\ast)=\mathbf 0$; conversely such a state does not move. (ii) Let $E=\tfrac12\mathbf v^{\top}\mathsf M\mathbf v+U(\mathbf x)$. Along trajectories $\dot E=\mathbf v\cdot(\mathbf F+\nabla U)=-\mathbf v^{\top}\mathsf D\,\mathbf v+\mathbf v\cdot\mathbf G\le0$ while $|\mathbf v|<\epsilon$. Near a strict minimum, small sublevel sets of $E$ around $(\mathbf x^\ast,\mathbf 0)$ trap trajectories.
\end{proof}

The hypotheses hold for SFM and SFM+TTC. The frictions of Eqs.~\eqref{eq:sfm}--\eqref{eq:wall} vanish at rest with power $-\kappa\,g\,[(\mathbf v_i-\mathbf v_j)\cdot\mathbf t_{ij}]^2\le0$ per pair and per wall, and the TTC force (Eq.~\eqref{eq:ttc}) is zero below relative speeds of $10^{-4}$\,m/s. With the densities frozen the SFM rest field is conservative,
\begin{equation}\label{eq:U}
U(\mathbf x)=\sum_i c_i\,|\mathbf x_i-\mathbf g|+\sum_{i<j}\phi_{ij}(d_{ij})+\sum_{i,W}\phi_{iW}(d_{iW}),\qquad c_i=\frac{m_iv_{{\rm eff},i}}{\tau_i},
\end{equation}
with $\phi$ the exponential and compression potentials ($\max|\mathbf F_0+\nabla U|\le\ResidCollMax$\,N over the \ResidBarN\ barrier and $\ResidPairMax$\,N over the \ResidPairN\ pair snapshots). SFM+TTC therefore has the fixed points and stability of SFM, so TTC and friction are at most class P. In force coupling ORCA's rest drive enters $\mathbf F_0$ and is a candidate for class L. Part (ii) does not cover ORCA's velocity-dependent constraint part.

\subsection{Effective Temperature of Per-Step Noise}\label{sec:teff}
Away from contact, one velocity component of agent $i$ under the update of Sect.~\ref{sec:model} is the autoregressive process $v_{n+1}=a\,v_n+\xi_n$, $a=1-\Delta t/\tau_i$, $\langle\xi_n\xi_{n'}\rangle=\sigma^2\delta_{nn'}$, with stationary variance $\sigma^2/(1-a^2)$. Setting $m_i\langle v^2\rangle=k_BT_{\rm eff}$ gives
\begin{equation}\label{eq:teff}
k_BT_{\rm eff}=\frac{m_i\,\sigma^2}{1-(1-\Delta t/\tau_i)^2}=\frac{m_i\,\sigma^2\tau_i}{2\,\Delta t}\left(1-\frac{\Delta t}{2\tau_i}\right)^{-1}.
\end{equation}
A restoring curvature $K$ adds terms of order $(\omega\Delta t)^2$, $\omega^2=K/m_i$, negligible in our snapshots. At fixed per-step $\sigma$, $k_BT_{\rm eff}\propto1/\Delta t$ up to the factor $(1-\Delta t/2\tau_i)^{-1}$; $\sqrt{\Delta t}$-scaled noise leaves only that factor. Eq.~\eqref{eq:teff} assumes frozen densities, linear friction and no contact or speed clipping. Through $\tau_i$, $T_{\rm eff}$ differs between agents by up to a factor of 1.5 (Table~\ref{tab:params}); $\sigma$ is set from the blocker's $\tau_i$. In separate runs the blocker's realised velocity variance is on average \TthreeRatioMean\ times Eq.~\eqref{eq:teff}; over the \TthreeCells\ noise settings of the escape experiments its median ratio is \TthreeCellsMedian. For activated escape, $\lambda\propto\exp(-\beta)$ with $\beta\equiv\Delta U/k_BT_{\rm eff}$, and the log rate ratio of two time steps has no free parameter,
\begin{equation}\label{eq:ell}
\ell\equiv\ln\frac{\lambda(\Delta t_1)}{\lambda(\Delta t_2)}=\Delta U\left(\frac{1}{k_BT_{\rm eff}(\Delta t_2)}-\frac{1}{k_BT_{\rm eff}(\Delta t_1)}\right),
\end{equation}
positive for $\Delta t_1<\Delta t_2$ at fixed per-step $\sigma$, near zero for $\sqrt{\Delta t}$-scaled noise, and absent for a deterministic resolver.

\subsection{Barriers from the String Method}\label{sec:barrier}
The string method~\cite{e2007} gives $\Delta U$ on the landscape of Eq.~\eqref{eq:U}, the SFM rest field, from the snapshot minimum to an exit, the blocker past the door plane, with all residual agents mobile (full), the neighbours pinned (pinned) or the blocker alone (alone). A string is used only if it converges (uphill perpendicular force below 5\,N, climbing image stationary at the maximum, under 1\% change on doubling the images), ends in an exit and, when alone, matches a direct two-dimensional path. Of the \FinBarLibN\ lone-trap snapshots sent to the string method, \FinBarGateFail\ failed this gate and are not used; \LoneBarN\ barriers enter the library. Kramers theory predicts~\cite{kramers1940,hanggi1990} exponential escape (Weibull shape 1)~\cite{weibull1951} and $\ln\lambda$ linear in $1/k_BT_{\rm eff}$ with slope $-\Delta U$ in each condition; the neighbours' configuration should put the full barrier below the alone one.
